\section{Paired histories and untuned head sensitivity}\label{app:matched}
The added experiment refits two median-loss LightGBM predictors, each with 300 trees, 31 leaves, learning rate .05, minimum leaf size 50, and seed 20260910. Both use exactly the same 600,000 sampled rows on days 1314--1554. One arm uses artificially censored sales for histories and fit targets; the other uses complete recorded sales. This removes artificial censoring, not possible natural stockouts in the original records. The two arms use the same 28 common feature definitions and six item-aggregate history features. Capacity, fill-ratio, censor-rate, and hit-rate features are excluded from both arms. Price features use the origin week. All history aggregates are rebuilt for each arm.

A disjoint 600,000-row sample on days 1555--1610 fits each arm's conditional absolute-error and exposure heads. Both use squared loss with either 220 or 660 trees; all other tree parameters remain fixed. The tree counts are not selected by validation or early stopping, so this comparison does not test an improvement in conditional-moment estimation. Their four combinations share the point predictor. Targets for these audit heads are complete sales and realized absolute error. Realized outcomes are noisy targets for conditional moments; substituting them directly into a test selector would not be a conditional-information oracle.

Calibration A/B and evaluation dates match the earlier development audit. All five rankings receive their largest complete-tie threshold meeting both calibration caps and the .35 case floor. The primary cap is .95 times the larger full-coverage A/B risk for that arm. The absolute grid is $\{.60,.62,.64012983268,.66,.68,.70\}$; relative factors are $\{.85,.90,.95,1,1.05\}$. Only the primary relative cap receives the full head factorial; all grid points use the 220-tree heads. Utility choices, model hashes, and thresholds are saved before each arm's new evaluation. The protocol fixes both arms before fitting. All dates were previously evaluated, so this is retrospective evidence, not a new external trial. Intervals use 2,000 paired item resamples, conditional on the fits, and are descriptive 95\% intervals without a simultaneous-risk claim.

\begin{table}[t]\centering\scriptsize
\caption{Untuned tree-count sensitivity at each arm's primary cap. Error/exposure head tree counts vary factorially, with the point predictor fixed; tree counts were not selected by validation or early stopping. Larger heads are not uniformly more accurate, so these comparisons do not test whether improved moment estimates resolve the choice. MSE is evaluated against realized error or sales, not observed conditional moments.}\label{tab:matched-heads}
\begin{tabular}{llrrlrr}\toprule
Arm & Trees $e/w$ & Error MSE & Exposure MSE & Case / exposure rule & $\Delta c$ & $\Delta d$\\\midrule
Artificially censored & 220 / 220 & 4.475 & 5.893 & Excess / Exposure descending & -39.95 & +7.05 \\
Artificially censored & 660 / 220 & 4.510 & 5.893 & Excess / Exposure descending & -40.66 & +8.29 \\
Artificially censored & 220 / 660 & 4.475 & 5.883 & Excess / Exposure descending & -37.94 & +8.08 \\
Artificially censored & 660 / 660 & 4.510 & 5.883 & Excess / Exposure descending & -40.04 & +8.14 \\
Complete recorded sales & 220 / 220 & 3.064 & 4.761 & Excess / Ratio & -15.81 & +9.72 \\
Complete recorded sales & 660 / 220 & 3.119 & 4.761 & Excess / Exposure descending & -34.52 & +11.45 \\
Complete recorded sales & 220 / 660 & 3.064 & 4.805 & Excess / Exposure descending & -34.85 & +9.88 \\
Complete recorded sales & 660 / 660 & 3.119 & 4.805 & Excess / Exposure descending & -34.98 & +10.71 \\\bottomrule\end{tabular}\end{table}

\begin{table}[t]\centering\scriptsize
\caption{All declared cap settings for the two refitted arms, using the 220-tree heads. Relative caps use the arm-specific maximum full A/B calibration risk. A missing contrast means no eligible policy, not equivalence.}\label{tab:matched-caps}
\begin{tabular}{llrlrr}\toprule
Arm & Cap specification & Risk cap & Case / exposure rule & $\Delta c$ & $\Delta d$\\\midrule
Artificially censored & absolute 0.600 & 0.6000 & Excess / Mixed & -12.25 & +9.86 \\
Artificially censored & absolute 0.620 & 0.6200 & Excess / Ratio & -28.79 & +10.45 \\
Artificially censored & absolute 0.640 & 0.6401 & Excess / Exposure descending & -40.23 & +10.83 \\
Artificially censored & absolute 0.660 & 0.6600 & Excess / Exposure descending & -40.37 & +7.70 \\
Artificially censored & absolute 0.680 & 0.6800 & Excess / Ratio & -18.64 & +4.15 \\
Artificially censored & absolute 0.700 & 0.7000 & Error / Error & +0.00 & +0.00 \\
Artificially censored & relative 0.850 & 0.5947 & Excess / Mixed & -12.19 & +10.26 \\
Artificially censored & relative 0.900 & 0.6297 & Excess / Ratio & -28.29 & +9.71 \\
Artificially censored & relative 0.950 & 0.6647 & Excess / Exposure descending & -39.95 & +7.05 \\
Artificially censored & relative 1.000 & 0.6997 & Error / Error & +0.00 & +0.00 \\
Artificially censored & relative 1.050 & 0.7347 & Error / Error & +0.00 & +0.00 \\
Complete recorded sales & absolute 0.600 & 0.6000 & Excess / Ratio & -22.87 & +14.92 \\
Complete recorded sales & absolute 0.620 & 0.6200 & Excess / Exposure descending & -36.26 & +12.36 \\
Complete recorded sales & absolute 0.640 & 0.6401 & Excess / Ratio & -13.21 & +8.75 \\
Complete recorded sales & absolute 0.660 & 0.6600 & Excess / Ratio & -4.06 & +4.41 \\
Complete recorded sales & absolute 0.680 & 0.6800 & Error / Error & +0.00 & +0.00 \\
Complete recorded sales & absolute 0.700 & 0.7000 & Error / Error & +0.00 & +0.00 \\
Complete recorded sales & relative 0.850 & 0.5661 & Excess / Ratio & -25.84 & +18.69 \\
Complete recorded sales & relative 0.900 & 0.5994 & Excess / Ratio & -23.03 & +14.98 \\
Complete recorded sales & relative 0.950 & 0.6327 & Excess / Ratio & -15.81 & +9.72 \\
Complete recorded sales & relative 1.000 & 0.6660 & Error / Error & +0.00 & +0.00 \\
Complete recorded sales & relative 1.050 & 0.6993 & Error / Error & +0.00 & +0.00 \\\bottomrule\end{tabular}\end{table}


The source package records ten model fits and zero external openings. The verifier independently reconstructs score masks, accepted masses and coverage differences, checks paired evaluation targets and training chronology, and replays the primary intervals. The post-fit moment audit additionally decomposes estimated-minus-realized excess into $(\hat E-E)-r(\hat M-M)$ on each selected set. This is descriptive accepted-mass accounting, not identification of the true conditional heads.

\paragraph{Accepted-mass diagnostics.}
For the complete-sales primary ratio policy, predicted accepted error exceeds realized error by 5.05\%, while predicted accepted exposure exceeds realized exposure by 3.32\%. The corresponding predicted and realized risks are .6390 and .6284. For descending exposure these biases are 6.52\% and 2.91\%, with risks .6516 and .6296. The aggregate discrepancy contains nonzero error and exposure terms; it cannot be described by its denominator term alone. These post-fit aggregate diagnostics do not identify conditional estimation error.

The complete-sales primary common set has 1,313,818 rows, exposure 2,225,652, and error 1,340,939.17. Case-only and exposure-only sets have exposure 135,897 and 419,613, respectively. Their union has risk .6467 versus cap .6327. Descriptive 95\% item-bootstrap risk intervals are [.6065,.6466] for the case policy, [.6106,.6482] for the exposure policy, and [.6283,.6668] for the union; all cross the cap. No joint population feasibility claim follows from the point-budget decomposition.
